\section*{INDEX OF TERMINOLOGY}
\phantomsection
\addcontentsline{toc}{section}{Index of Terminology}

The reference numbers indicate the chapter, section, and sub-section (or exercise) in that order. (R refers to Summary of Results.)

\begin{Shaded}
\begin{Highlighting}[]
\NormalTok{Adjunction of a greatest element to an ordered set : III.1.7}
\NormalTok{Adjunction (of one set to another) : R.5.4}
\NormalTok{Agreeing on a set (functions) : II.3.5}
\NormalTok{Agreement (of two functions on a set) : R.2.14}
\NormalTok{Aleph : III.6. Ex. 10}
\NormalTok{Antecedent assemblies : I. App. 4}
\NormalTok{Antidirected (ordered set) : III.1. Ex. 23}
\NormalTok{Applying the results of one theory in another : I.2.4}
\NormalTok{Argument : R.1.2}
\NormalTok{Assembly : I.1.1}
\NormalTok{antecedent : I. App. 4}
\NormalTok{balanced : I. App. 4}
\NormalTok{of the first (second) species : I.1.3}
\NormalTok{perfectly balanced : I. App. 4}
\NormalTok{Associated equivalence relation : II.6.2}
\NormalTok{Associativity criterion for product structures : IV.2.4}
\NormalTok{of the intersection of a family of sets : R.4.8}
\NormalTok{of the product of a family of sets : R.4.11}
\NormalTok{of the union and intersection of two sets : R.1.14}
\NormalTok{of the union of a family of sets : R.4.3}
\NormalTok{Automorphism : IV.1.5, R.8.6}
\NormalTok{Auxiliary base sets : IV.1.3,1.4}
\NormalTok{constant, method of : I.3.3}
\NormalTok{hypothesis, method of : I.3.3}
\NormalTok{Axiom, explicit : I.2.1}
\NormalTok{implicit : I.2.1}
\NormalTok{of choice : R.4.10}
\NormalTok{of extent : II.1.3}
\NormalTok{of infinity : III.6.1}
\end{Highlighting}
\end{Shaded}

\begin{Shaded}
\begin{Highlighting}[]
\NormalTok{Axiom of the ordered pair : II.2.1}
\NormalTok{of the set of subsets : II.5.1}
\NormalTok{of the set of two elements : II.1.5}
\NormalTok{Axioms, equivalent : R.8.4}
\NormalTok{of structures of the same species : R.8.2}
\NormalTok{Axiom(s) of a species of structures : IV.1.4}

\NormalTok{Balanced assembly : I. App. 4}
\NormalTok{word : I. App. 2}
\NormalTok{Base of an expansion, or of a system of numeration : III.5.7}
\NormalTok{of a scale of sets : R.8.1}
\NormalTok{sets (of an echelon construction) : IV.1.1}
\NormalTok{sets, auxiliary : IV.1.3, 1.4}
\NormalTok{sets, principal : IV.1.3, 1.4}
\NormalTok{Bijection : II.3.7, R.2.9}
\NormalTok{Bijective mapping : II.3.7, R.2.9}
\NormalTok{Binomial coefficient : III.5.8}
\NormalTok{Boolean lattice : III.1. Ex. 17}
\NormalTok{Bound, greatest lower and least upper (of a set or a mapping) : III.1.9, R.6.7}
\NormalTok{least upper (of a set of cardinals) : III.3.2}
\NormalTok{lower : III.1.8, R.6.7}
\NormalTok{strict upper : III.2.4}
\NormalTok{upper : III.1.8, R.6.7}
\NormalTok{Bounded, bounded above, bounded below (set or mapping) : III.1.8, R.6.7}
\NormalTok{Branched (ordered set) : III.1. Ex. 24}
\end{Highlighting}
\end{Shaded}

Canonical decomposition of a function : II.6.5, R.5.3 Canonical extension, of a correspondence to sets of subsets : II.5.1 of a family of functions to the product sets : II.5.7 of mappings : IV.1.2 of two functions to the product sets : II.3.9 signed : IV.2. Ex. 1 Canonical injection : II.3.7 Canonical mapping, IV, App. 4 of a subset of E into E : II.3.7, R.2.3 of A\(^{B \times C}\) onto (A\(^{B}\))\(^{C}\) : II.5.2, R.4.14 of G onto G (G a graph) : II.3.7 of F\(^{E}\) onto F(E, F) : II.5.2 of F(B × C, A) onto F(B, F(C, A)) : II.5.2 of ∏\emph{\{i∈\{α\}\} X\_i onto X\_α : II.5.3 of ∏}\{i∈\{α, β\}\} X\_i onto X\_α × X\_β : II.5.3

Canonical mapping of \(\prod_{i\in I}X_t\) onto \(\prod_{\lambda \in L}\left(\prod_{i\in J_\lambda}X_t\right)\): II.5.5, R.4.11\\
of \(\prod_{i\in I}X_t\) onto \(\left(\prod_{i\in J_\alpha}X_t\right) \times \left(\prod_{i\in J_\beta}X_t\right)\) ((J\(_\alpha\), J\(_\beta\)) a partition of I): II.5.5\\
of \(\prod_{i\in \{\alpha,\beta,\gamma\}}X_t\) onto \(X_\alpha \times X_\beta \times X_\gamma\): II.5.5\\
of \(\prod_{i\in I}X_t^E\) onto \(\left(\prod_{i\in I}X_t\right)^E\): II.5.7, R.4.13\\
of E onto E/R: II.6.2, R.5.2\\
of A/R\(_A\) onto f\textless A\textgreater: II.6.6, R.5.5\\
of (E/S)/(R/S) onto E/R: II.6.7, R.5.9\\
of (E × E')/(R × R') onto (E/R) × (E'/R'): II.6.8, R.5.10\\
into a direct limit of a direct limit obtained by restriction of the index set: III.7.6\\
of E\(_\beta\) into lim E\(_\alpha\): III.7.5, R.6.13\\
of an inverse limit into an inverse limit obtained by restriction of the index set: III.7.1\\
of lim E\(_\alpha\) into E\(_\beta\): III.7.1, R.6.14\\
of E × F onto F × E: R.3.4\\
of E × F × G onto (E × F) × G, etc.: R.3.12\\
of FA onto \(\prod_{i\in I}FA_i\) (where A = \(\bigcup_{i\in I}A_i\), and A\(_t\) ∩ Ax = φ whenever t ≠ x): R.4.15\\
symmetry: R.3.4\\
Cantor's theorem: III.3.6\\
Cardinal, dominant: III.6.Ex.21\\
finite: III.4.1\\
inaccessible: III.6.Ex.22\\
product: III.3.3\\
regular: III.6.Ex.17\\
singular: III.6.Ex.17\\
strongly inaccessible: III.6.Ex.22\\
Cardinal of a set: III.3.1\\
Cardinal sum: III.3.3\\
Chain of an element in an ordered set (with respect to a mapping): III.2.Ex.6\\
Characteristic function of a subset of a set: III.5.5\\
Characterization, typical: IV.1.4\\
Choice, axiom of: R.4.9\\
Class, equivalence: II.6.2, R.5.2\\
of objects equivalent to x (with respect to an equivalence relation): II.6.9\\
Closed interval: III.1.13, R.6.4\\
Closure (mapping of an ordered set into itself): III.1.Ex.13

Coarser equivalence relation : II.6.7 preordering : III.1.4 structure : IV.2.2 Coefficient, binomial : III.5.8 Cofinal subset : III.1.7, R.6.5 Coincidence (of two functions on a set) : II.3.5, R.2.14 Coinitial subset : III.1.7, R.6.5 Collectivizing relation : II.1.4 Commutativity (of union and intersection) : R.1.14 Comparable elements : III.1.12 structures : IV.2.2 Compatibility, of a function with two equivalence relations : II.6.5, R.5.8 of a relation with an equivalence relation : II.6.5, R.5.7 Complement of a set : II.1.7, R.1.7 Complete lattice : III.1.Ex.11 Complete solution : I.5.2 Completely ramified (ordered set) : III.2.Ex.8 Completion of an ordered set : III.1.Ex.15 Composition, of mappings : R.2.11 of sets : R.3.10 of two correspondences : II.3.3 of two graphs : II.3.3 Conjunction of two relations : I.3.4 Connected components of a set with respect to a relation : II.6.Ex.10 Constant, auxiliary : I.3.3 function or mapping : II.3.4, R.2.3 of a theory : I.2.1 Constituents of a set with respect to a relation : II.6.Ex.11 Construction, echelon : IV.1.1 formative : I.1.3 Contained in a set : II.1.2, R.1.12 Continuum hypothesis : III.6.4 Continuum, power of : III.6.4 Contradictory axioms : R.8.6 theory : I.2.2 Contravariant signed echelon type : IV.2.Ex.1 Coordinate (first, second) of an ordered pair : II.2.1 Coordinate function (first, second) : II.3.6 of index i : II.5.3 Coordinate functions, on a product of a family of sets : R.4.11 on a product of several sets : R.3.12 on a product of two sets : R.3.1 Corollary : I.2.2 Correspondence, between two sets : II.3.1

Correspondence defined at an object x : II.3.1 defined by a relation : II.3.1 inverse : II.3.2 one-to-one : II.3.7, R.2.9 Correspondences, composition of : II.3.3 Countable set : III.6.4, R.7.7 Countable union, intersection, product : R.7.9 Covariant signed echelon type : IV.2.Ex.1 Covering : II.4.6, R.4.4 finer : II.4.6 of a set : II.4.6 Criterion, deductive : I.2.2 formative : I.1.4 of deduction : I.3.3 of substitution : I.1.2 Critical ordinal : III.6.Ex.13

Decent set : III.2.Ex.20 Decimal system : III.5.7 Decomposition, canonical : II.6.5, R.5.3 Decreasing family of subsets : III.1.5, R.6.12 Decreasing mapping : III.1.5, R.6.12 Deduced (structure) : IV.1.6 Deduction, criterion of : I.3.3 procedure : IV.1.6 Deductive criterion : I.2.2 Definition : I.1.1 Degree of disjointness of a covering : III.6.Ex.25 Demonstrative text : I.2.2 Descending induction : III.4.3 Diagonal, of A × A : II.3.3, R.3.4 of E¹ : II.5.3 Diagonal mapping, of A into A × A : II.3.7, R.3.4 of E into E¹ : II.5.3 Diagrams : R.2.2 Difference of two integers : III.5.2 Different from : I.5.1 Digit : III.5.7 Direct image of a structure : IV.2.6 Direct limit, of a direct system of mappings : III.7.6, R.6.13 of a direct system of sets : III.7.5, R.6.13 Direct system, of mappings : III.7.6, R.6.13 of sets : III.7.5, R.6.13 of subsets : III.7.6

Directed (left, right) : III.1.10, R.6.8 with respect to the relation ≤ : III.1.10 Disjoint segments : I.App.1 sets : II.4.7, R.1.13 Disjunction, of cases, method of : I.3.3 of two relations : I.1.3 Distributive lattice : III.1.Ex.16 Distributivity, of union and intersection of a family of sets : R.4.3, 4.8 of union and intersection of two sets : R.1.14 Divergent mapping (of an ordinal into itself) : III.6.Ex.23 Divisible by an integer : III.5.6 Division, Euclidean : III.5.6 of an integer : III.5.6 Domain, of a correspondence : II.3.1 of a graph : II.3.1 Dominant cardinal : III.6.Ex.21 Double family : II.3.4 Double sequence : III.6.1, R.7.8 Dual formulae : R.1.15, R.4.7 Duality rule : R.1.15, R.4.7 Dyadic system : III.5.7

Echelon, echelon construction, echelon construction scheme : IV.1.1 signed : IV.2.Ex.1 type : IV.1.Ex.1 Element : II.1.1 fixed (under a mapping or set of mappings) : R.2.3 generic : R.1.2 greatest : III.1.3, R.6.5 invariant (under a mapping or set of mappings) : R.2.3 least : III.1.3, R.6.5 maximal : III.1.6, R.6.6 minimal : III.1.6, R.6.6 Elements, comparable : III.1.12 Empty function : II.3.4 Empty set : II.1.7, R.1.8 Empty word : I.App.1 Endowed with a structure : IV.1.4 Endpoints of an interval : III.1.13, R.6.4 Equal : I.5.1 Equalitarian theory : I.5.1 Equality, relation of : I.5.1, R.1.6 Equation : I.5.2 Equipotent sets : III.3.1, R.7.1

\begin{Shaded}
\begin{Highlighting}[]
\NormalTok{Equivalence, on a set: II.6.1}
\NormalTok{Equivalence class: II.6.2, R.5.2}
\NormalTok{Equivalence relation: II.6.1, R.5.2}
\NormalTok{associated with a function: II.6.2}
\NormalTok{coarser: II.6.7}
\NormalTok{finer: II.6.7}
\NormalTok{induced: II.6.6}
\NormalTok{on a set: II.6.1}
\NormalTok{product: II.6.8}
\NormalTok{quotient: II.6.7}
\NormalTok{Equivalent axioms: R.8.4}
\NormalTok{Equivalent elements: II.6.1}
\NormalTok{Equivalent powers: R.7.2}
\NormalTok{Equivalent relations: I.3.5, R.1.3}
\NormalTok{Equivalent species of structures: IV.1.7}
\NormalTok{Equivalent theories: I.2.4}
\NormalTok{Euclidean division: III.5.6}
\NormalTok{Even integer: III.5.6}
\NormalTok{Existential quantifier: I.4.1}
\NormalTok{Expansion of an integer to base b: III.5.7}
\NormalTok{Explicit axiom: I.2.1}
\NormalTok{Exponentiation: R.4.9}
\NormalTok{Extension, canonical: IV.1.2}
\NormalTok{inverse (of a mapping to sets of subsets): R.2.6}
\NormalTok{of a correspondence or mapping to sets of subsets: II.5.1, R.2.4}
\NormalTok{of a family of functions to product sets: II.5.7}
\NormalTok{of a mapping: II.3.5, R.2.13}
\NormalTok{of an ordering, order relation, preordering, preorder relation: III.1.4, R.6.1}
\NormalTok{of several mappings to the product sets: R.3.14}
\NormalTok{of two functions to product sets: II.3.9}
\NormalTok{Extent, axiom of: II.1.3}

\NormalTok{Factor, of an integer: III.5.6}
\NormalTok{of a product: II.2.2 and II.5.3}
\NormalTok{Factors, of a product set: R.3.1, 4.9}
\NormalTok{Factorial n: III.5.8}
\NormalTok{Factorization, canonical: R.5.3}
\NormalTok{of a mapping: R.2.11}
\NormalTok{False relation: I.2.2}
\NormalTok{Family: II.3.4}
\NormalTok{double: II.3.4}
\NormalTok{finite: III.4.1}
\NormalTok{of elements of a set: II.3.4, R.2.14}
\NormalTok{of mutually disjoint sets: II.4.7}
\end{Highlighting}
\end{Shaded}

\begin{Shaded}
\begin{Highlighting}[]
\NormalTok{Family of subsets, decreasing (increasing) : III.1.5}
\NormalTok{of subsets of a set : II.3.4}
\NormalTok{Final character (of a totally ordered set) : III.6.Ex.16}
\NormalTok{Final segment of a word : I.App.1}
\NormalTok{Final structure : IV.2.5}
\NormalTok{Finer covering : II.4.6}
\NormalTok{Finer preordering : III.1.4}
\NormalTok{Finer relation : II.6.7}
\NormalTok{Finer structure : IV.2.2}
\NormalTok{Finite cardinal : III.4.1}
\NormalTok{Finite character (property of, set of subsets of) : III.4.5, R.6.11}
\NormalTok{Finite family : III.4.1}
\NormalTok{Finite sequence, finite sequence of elements of a set : III.5.4, R.7.8}
\NormalTok{Finite set : IV.4.1}
\NormalTok{First coordinate of an ordered pair : II.2.1}
\NormalTok{First projection of a graph : II.3.1}
\NormalTok{First species : I.1.3}
\NormalTok{First term of a sequence : III.5.4}
\NormalTok{Fixed element : II.3.4, R.2.3}
\NormalTok{Formative construction : I.1.3}
\NormalTok{criterion : I.1.4}
\NormalTok{Free subset of an ordered set : III.1.Ex.5}
\NormalTok{Function : see mapping}
\NormalTok{characteristic : III.5.5}
\NormalTok{coordinate : II.3.6, II.5.3}
\NormalTok{defined on A with values in B : II.3.4}
\NormalTok{not depending on x : II.3.9}
\NormalTok{of several variables : R.3.13}
\NormalTok{of two arguments : II.3.9}
\NormalTok{Functional graph : II.3.4}
\NormalTok{Functional relation : I.5.3, R.2.1}
\NormalTok{Functional symbol : I.5.3}

\NormalTok{General solution : I.5.2}
\NormalTok{General term (of a sequence) : R.7.8}
\NormalTok{Generalized continuum hypothesis : III.6.4}
\NormalTok{Generic element (of a set) : R.1.2}
\NormalTok{Generic structure : IV.1.4}
\NormalTok{Graph : II.3.1}
\NormalTok{functional : II.3.4}
\NormalTok{inverse : II.3.2}
\NormalTok{of a correspondence : II.3.1}
\NormalTok{of a mapping : R.3.5}
\NormalTok{of a relation : II.3.1, R.3.2}
\end{Highlighting}
\end{Shaded}

Graph symmetric : II.3.2 Graphs, composition of : II.3.3 Greater : III.1.3, R.6.3 Greatest element of an ordered set : III.1.7, R.6.5 Greatest lower bound : III.1.9, R.6.7

Half-open interval : III.1.13, R.6.4 Hypothesis, auxiliary : I.3.3 continuum : III.6.4 generalized continuum : III.6.4 inductive : III.2.2, III.4.3

Identification : R.8.5 Identity : R.1.3 Identity correspondence : II.3.3 Identity mapping : II.3.4, R.2.3 Image, direct (of a structure) : IV.2.6 inverse (of a set under a mapping) : R.2.6 inverse (of a structure) : IV.2.4 of a covering : II.4.6 of an equivalence relation : II.6.6 of a function (by abuse of language) : R.2.4 of a set : II.3.2 of a set under a correspondence : II.3.1 of a set under a function : II.3.4, R.2.4 of a set under a graph : II.3.1 Implicit axiom : I.2.1 Imply : R.1.3 Inaccessible cardinal : III.6.Ex.22 Inaccessible initial ordinal : III.6.Ex.16 Inclusion relation : II.1.2, III.1.1, R.1.2 Inconsistent axioms : R.8.6 Increasing family of subsets : III.1.5, R.6.12 Increasing mapping : III.1.5, R.6.12 Indecomposable ordinal : III.2.Ex.16 Independent (of other axioms) : I.2.Ex.1 Index set of a family : II.3.4, R.2.2 Indicial notation : R.2.2 Induced equivalence relation : R.5.5 Induced mapping : II.6.5 Induced ordering, order relation, preordering, preorder relation: III.1.4, R.6.1 Induced relation : II.6.3 Induced structure : IV.2.4

Induction, descending : III.4.3 Noetherian : III.6.5 principle of : III.2.2, III.4.3 restricted : III.4.3 starting at k : III.4.3 transfinite : III.2.2 Inductive hypothesis : III.2.2, III.4.3 Inductive set : III.2.4, R.6.9 Infimum : III.1.9, R.6.7 Infinite sequence : III.6.1, R.7.8 Infinite set : III.6.1 Infinity, axiom of : III.6.1 Initial ordinal : III.6.Ex.10 Initial segment of a word : I.App.1 Initial structure : IV.2.3 Injection : II.3.7, R.2.8 canonical : II.3.7 Injective mapping : II.3.7, R.2.8 Integer, natural : III.4.1 Integral part of the quotient of two integers : III.5.6 Intersecting subsets : R.1.13 Intersection, countable : R.7.9 of a family of sets, or subsets : II.4.1, R.4.6 of a set of sets : II.4.1 of several sets : R.1.13 Intervals (closed, half-open, open, unbounded) : III.1.13, R.6.4 Intransitive: II.6.Ex.11 Intrinsic term : IV.1.6 Invariant element : R.2.3 Inverse extension of a mapping : R.2.6 Inverse, of a (bijective) mapping : II.3.7, R.2.9 of a correspondence : II.3.2 of a graph : II.3.2 Inverse image : R.2.6 of a covering : II.4.6 of an equivalence relation : II.6.6 of a set : II.3.2 of a structure : IV.2.4 Inverse isomorphisms : IV.1.5 Inverse limit, of a family of mappings : III.7.2, R.6.14 of a family of sets : III.7.1, R.6.14 Inverse system, of mappings : III.7.2, R.6.14 of sets : III.7.1, R.6.14 of subsets : III.7.2

Involutory permutation : II.3.7, R.2.9 Irreducible element (of a lattice) : III.4.Ex.7 Isomorphic sets : IV.1.5 structures : IV.1.5, R.8.6 Isomorphism : IV.1.5, R.8.6 of ordered sets : III.1.3 Iterates of a mapping : III.6.2, R.2.11

Larger than : III.1.3 Last term of a finite sequence : III.5.4 Lattice : III.1.11, R.6.8 Boolean : III.1.Ex.17 complete : III.1.Ex.11 distributive : III.1.Ex.16 relatively complemented : III.1.Ex.17 Least element of an ordered set : III.1.7, R.6.5 Least upper bound : III.1.9, R.6.7 of a family of cardinals : III.3.2 Left directed set : III.1.10, R.6.8 Left-hand endpoint of an interval : III.1.13 Left inverse (of an injection) : II.3.8 Legitimation, theorem of : I.3.3 Lemma : I.2.2 Zorn's : III.2.4, R.6.10 Length, of a finite sequence : III.5.4 of a word : I.App.1 Less than : III.1.3, R.6.3 Lexicographical order relation, ordering : III.2.6 product : III.2.6 and III.2.Ex.10 Limit, direct : III.7.5, R.6.13 inverse : III.7.1, R.6.14 Linearly ordered : see totally ordered Link : I.1.1 Logical components of a relation : I.App.Ex.6 Logical construction : I.App.Ex.6 Logical sign : I.1.1 Logical theory : I.3.1 Logically constructed : I.App.Ex.6 Logically irreducible : I.App.Ex.6 Lower bound : III.1.8, R.6.7 greatest : III.1.9, R.6.7

Mapping : R.2.1 bijective : II.3.7, R.2.9

\begin{Shaded}
\begin{Highlighting}[]
\NormalTok{bounded above (bounded below, bounded) : III.1.8, R.6.7}
\NormalTok{compatible with an equivalence relation : II.6.5}
\NormalTok{compatible with two equivalence relations : R.5.8}
\NormalTok{composite : II.3.7, R.2.11}
\NormalTok{constant : II.3.4, R.2.3}
\NormalTok{decreasing : III.1.5, R.6.12}
\NormalTok{diagonal : II.3.7, II.5.3, R.3.4}
\NormalTok{empty : II.3.4}
\NormalTok{identity : II.3.4, R.2.3}
\NormalTok{increasing : III.1.5, R.6.12}
\NormalTok{injective : II.3.7, R.2.8}
\NormalTok{inverse : II.3.7, R.2.9}
\NormalTok{monotone : III.1.5}
\NormalTok{of a set into a set : II.3.4}
\NormalTok{of a set onto a set : II.3.7}
\NormalTok{one{-}to{-}one : II.3.7, R.2.9}
\NormalTok{order{-}preserving : III.1.5}
\NormalTok{order{-}reversing : III.1.5}
\NormalTok{partial : II.3.9, R.3.13}
\NormalTok{strictly decreasing (strictly increasing, strictly monotone) : III.1.5, R.6.12}
\NormalTok{surjective : II.3.7, R.2.4}
\NormalTok{universal : IV.3.1}
\NormalTok{Mappings, agreeing on a set : II.3.5}
\NormalTok{canonical : see canonical}
\NormalTok{Mathematical theory : I.1.1, I.2.1, I.2.2}
\NormalTok{Maximal element : II.1.6, R.6.7}
\NormalTok{Maximum : R.6.5}
\NormalTok{Membership, relation of : II.1.1, R.1.10}
\NormalTok{Method, of disjunction of cases : I.3.3}
\NormalTok{of reductio ad absurdum : I.3.3}
\NormalTok{of the auxiliary constant : I.3.3}
\NormalTok{of the auxiliary hypothesis : I.3.3}
\NormalTok{Minimal element : III.1.6, R.6.6}
\NormalTok{Minimum : R.6.5}
\NormalTok{Mobile (set of finite subsets) : III.4.Ex.11}
\NormalTok{Model of a theory : I.2.4}
\NormalTok{Monotone mapping : III.1.5}
\NormalTok{Morphism : IV.2.1}
\NormalTok{Multiple of an integer : III.5.6}
\NormalTok{Multiple sequence : III.6.1}
\NormalTok{Multivalent theory : R.8.7}
\NormalTok{Mutually disjoint (family of sets) : II.4.7}
\NormalTok{Natural integer : III.4.1}
\end{Highlighting}
\end{Shaded}

\begin{Shaded}
\begin{Highlighting}[]
\NormalTok{Negation of a relation : I.1.3}
\NormalTok{Noetherian induction, principle of : III.6.5}
\NormalTok{Noetherian set : III.6.5}
\NormalTok{nth iterate of a mapping : III.6.2, R.2.11}
\NormalTok{nth term of a sequence : III.6.1, R.7.8}
\NormalTok{Number of elements of a finite set : III.4.1}
\NormalTok{Numerical symbol : III.5.7}
\NormalTok{Numeration, system of : III.5.7}

\NormalTok{Odd integer : III.5.6}
\NormalTok{One{-}to{-}one correspondence, mapping : II.3.7, R.2.9}
\NormalTok{Open interval : III.1.13, R.6.4}
\NormalTok{Opposite order relation, preorder relation : III.1.1, III.1.2, R.6.1}
\NormalTok{Ordered pair : II.2.1}
\NormalTok{Order{-}preserving mapping : III.1.5}
\NormalTok{Order relation, associated with a preorder relation : III.1.2}
\NormalTok{between x and y, with respect to x and y : III.1.1}
\NormalTok{induced on a set : III.1.4}
\NormalTok{lexicographical : III.2.6}
\NormalTok{on a set : III.1.1, R.6.1}
\NormalTok{opposite to an order relation : III.1.1}
\NormalTok{Order{-}reversing mapping : III.1.5}
\NormalTok{Order{-}structure : R.6.1}
\NormalTok{total : III.1.12}
\NormalTok{Order type : III.2.Ex.13}
\NormalTok{Ordered set : III.1.3, R.6.1}
\NormalTok{antidirected : III.1.Ex.23}
\NormalTok{branched : III.1.Ex.24}
\NormalTok{completely ramified : III.2.Ex.8}
\NormalTok{partially well{-}ordered : III.2.Ex.4}
\NormalTok{ramified : III.2.Ex.8}
\NormalTok{scattered : III.1.Ex.20}
\NormalTok{without gaps : III.1.Ex.19}
\NormalTok{Ordering : III.1.1, R.6.1}
\NormalTok{induced : III.1.4}
\NormalTok{lexicographical : III.2.6}
\NormalTok{product : III.1.4}
\NormalTok{total : III.1.12}
\NormalTok{Ordinal : III.2.Ex.14}
\NormalTok{critical : III.6.Ex.13}
\NormalTok{functional symbol : III.2.Ex.17}
\NormalTok{inaccessible initial : III.6.Ex.16}
\NormalTok{indecomposable : III.2.Ex.16}
\NormalTok{Ordinal, initial : III.6.Ex.10}
\end{Highlighting}
\end{Shaded}

Ordinal product (of order-types) : III.2.Ex.13 regular : III.6.Ex.16 singular : III.6.Ex.16 Ordinal sum, of a family of non-empty ordered sets indexed by an ordered set : III.1.Ex.3 (of order-types) : III.2.Ex.13

\begin{Shaded}
\begin{Highlighting}[]
\NormalTok{Pair, ordered : II.2.1}
\NormalTok{Pairwise disjoint (family of subsets) : R.4.4}
\NormalTok{Parameter, parametric set, parametric representation : II.3.7, R.2.14}
\NormalTok{Partial mapping : II.3.9, R.3.13}
\NormalTok{Partial product : II.5.4}
\NormalTok{Partially well{-}ordered (ordered set) : III.2.Ex.1}
\NormalTok{Partition of a set : II.4.7, R.4.4}
\NormalTok{Passage to quotient sets : II.6.3, II.6.5, R.5.7, R5.8}
\NormalTok{Perfectly balanced assembly : I.App.4}
\NormalTok{Permutation : II.3.7, R.2.9}
\NormalTok{Poorer (species of structures) : IV.1.6}
\NormalTok{Power, of a set : III.3.1, R.7.2}
\NormalTok{of the continuum : III.6.4}
\NormalTok{Powers, equivalent : R.7.2}
\NormalTok{sum of : R.7.5}
\NormalTok{Predecessor (of an ordinal) : III.2.Ex.14}
\NormalTok{Preorder relation : III.1.2, R.6.1}
\NormalTok{on a set : III.1.2}
\NormalTok{opposite to a preorder relation : III.1.2}
\NormalTok{Preordered set : R.6.1}
\NormalTok{Preordering : III.1.2}
\NormalTok{(coarser, finer) : III.1.4}
\NormalTok{Principal base sets : IV.1.3, 1.4}
\NormalTok{Principle of induction : III.4.3}
\NormalTok{of Noetherian induction : III.6.5}
\NormalTok{of transfinite induction : III.2.2}
\NormalTok{Procedure of deduction of structures : IV.1.6}
\NormalTok{Product cardinal : III.3.3}
\NormalTok{Product, lexicographic : III.2.6}
\NormalTok{Product of a family of mappings : II.5.7}
\NormalTok{of a family of sets : II.5.3, R.4.9}
\NormalTok{of cardinals : III.3.3}
\NormalTok{of order relations, orderings, preorder relations, preorderings : III.1.4}
\NormalTok{of ordered sets, preordered sets : III.1.4}
\NormalTok{of several sets : R.3.12}
\NormalTok{of two coverings : II.4.6}
\NormalTok{of two equivalence relations : II.6.8, R.5.10}
\end{Highlighting}
\end{Shaded}

\begin{Shaded}
\begin{Highlighting}[]
\NormalTok{Product of two sets : II.2.2, R.3.1}
\NormalTok{Product ordinal : III.2.Ex.13}
\NormalTok{Product structure, product of structures : IV.2.4}
\NormalTok{Product partial : II.5.4}
\NormalTok{Projection, (first, second) of a graph : II.3.1}
\NormalTok{(first, second) of an ordered pair : II.2.1}
\NormalTok{onto a factor : II.5.3}
\NormalTok{onto a partial product : II.5.4}
\NormalTok{Projections : R.3.1, R.3.12, R.4.11}
\NormalTok{Proof : I.2.2}
\NormalTok{Proper segment : I.App.1}
\NormalTok{Property of finite character : III.4.5}
\NormalTok{Proposition : I.2.2}
\NormalTok{Pseudo{-}ordinal : III.2.Ex.20}

\NormalTok{Quantified theory : I.4.2}
\NormalTok{Quantifier, existential : I.4.1}
\NormalTok{typical : I.4.4}
\NormalTok{universal : I.4.1}
\NormalTok{Quotient equivalence relation : II.6.7, R.5.9}
\NormalTok{Quotient of a preorder relation or preordered set by an equivalence relation : III.1.Ex.2}
\NormalTok{Quotient of two integers : III.5.6}
\NormalTok{Quotient set : II.6.2, R.5.2}
\NormalTok{Quotient structure : IV.2.6}

\NormalTok{Ramified (ordered set) : III.2.Ex.8}
\NormalTok{Range, of a correspondence : II.3.1}
\NormalTok{of a graph : II.3.1}
\NormalTok{Realization, of an echelon type : IV.1.Ex.1}
\NormalTok{of a signed echelon type : IV.2.Ex.1}
\NormalTok{Reductio ad absurdum, method of : I.3.3}
\NormalTok{Refinement of a covering : II.4.6}
\NormalTok{Reflexive relation on a set : II.6.1, R.5.1}
\NormalTok{Regular cardinal : III.6.Ex.17}
\NormalTok{Regular ordinal : IV.6.Ex.16}
\NormalTok{Relation : I.1.3}
\NormalTok{between an element of A and an element of B : II.3.1}
\NormalTok{collectivizing : II.1.4}
\NormalTok{compatible with an equivalence relation : II.6.3, R.5.7}
\NormalTok{compatible with two equivalence relations : II.6.8}
\NormalTok{equivalence : II.6.1, R.5.2}
\NormalTok{false : I.2.2}
\NormalTok{functional : I.5.3, R.2.1}
\end{Highlighting}
\end{Shaded}

\begin{Shaded}
\begin{Highlighting}[]
\NormalTok{Relation, having a graph : II.3.1}
\NormalTok{    induced by an equivalence relation on passing to the quotient : II.6.3}
\NormalTok{    of equality : I.5.1, R.1.6}
\NormalTok{    of inclusion : II.1.2, III.1.1, R.1.12}
\NormalTok{    of membership : II.1.1, R.1.10}
\NormalTok{    order : III.1.1, R.6.1}
\NormalTok{    preorder : III.1.2, R.6.1}
\NormalTok{    reflexive : II.6.1, R.5.1}
\NormalTok{    single{-}valued : I.5.3}
\NormalTok{    symmetric : II.6.1, R.3.4}
\NormalTok{    total order : II.1.12}
\NormalTok{    transitive : II.6.1, R.5.1}
\NormalTok{    transportable : IV.1.3}
\NormalTok{    true : I.2.2}
\NormalTok{    well{-}ordering : III.2.1}
\NormalTok{    ≤ (resp. ≥, ⊂, ⊃) : III.1.3}
\NormalTok{Relations, equivalent : I.3.5, R.1.3}
\NormalTok{Relational sign : I.1.3}
\NormalTok{Relative complement (of an element in a lattice) : III.1.Ex.17}
\NormalTok{Relatively complemented lattice : III.1.Ex.17}
\NormalTok{Remainder (on division of a by b) : III.5.6}
\NormalTok{Representation, parametric : II.3.7, R.2.14}
\NormalTok{Representative, of an equivalence class : II.6.2}
\NormalTok{Representatives, system of : II.6.2}
\NormalTok{Restricted induction : III.4.3}
\NormalTok{Restriction, of a direct system : III.7.6}
\NormalTok{    of a function : II.3.5, R.2.13}
\NormalTok{    of an inverse system : III.7.1}
\NormalTok{Retraction of an injection : II.3.8}
\NormalTok{Richer structure : R.8.3}
\NormalTok{Right directed set : III.1.10, R.6.8}
\NormalTok{Right{-}hand endpoint of an interval : III.1.13}
\NormalTok{Right inverse (of a surjection) : II.3.8}

\NormalTok{Saturated set : R.5.6}
\NormalTok{Saturation of a set with respect to an equivalence relation : R.5.6}
\NormalTok{Scale of sets : R.8.1}
\NormalTok{Scattered (ordered set) : III.1.Ex.20}
\NormalTok{Scheme : I.2.1}
\NormalTok{    echelon construction : IV.1.1}
\NormalTok{    of selection and union : II.1.6}
\NormalTok{Second coordinate of an ordered pair : II.2.1}
\NormalTok{    projection : II.3.1}
\NormalTok{    species : I.1.3}
\end{Highlighting}
\end{Shaded}

\begin{Shaded}
\begin{Highlighting}[]
\NormalTok{Section, of a correspondence : II.3.1}
\NormalTok{of a graph : II.3.1}
\NormalTok{of a set (with respect to an equivalence relation) : II.6.2}
\NormalTok{of a subset of a product : R.3.7}
\NormalTok{of a surjection : II.3.8}
\NormalTok{Segment, final : I.App.1}
\NormalTok{initial : I.App.1}
\NormalTok{of a word : I.App.1}
\NormalTok{of an assembly : I.App.4}
\NormalTok{proper : I.App.1}
\NormalTok{with endpoint x : III.2.1}
\NormalTok{Segments, disjoint : I.App.1}
\NormalTok{Selection and union, scheme of : II.1.6}
\NormalTok{Separation (of the elements of E by α{-}mappings) : IV.3.3}
\NormalTok{Sequence : R.7.8}
\NormalTok{double : R.7.8}
\NormalTok{finite : III.5.4, R.7.8}
\NormalTok{infinite : III.6.1, R.7.8}
\NormalTok{obtained by ranging a countable family in an order defined by a mapping : III.6.1}
\NormalTok{of elements of a set : III.6.1}
\NormalTok{significant : I.App.2}
\NormalTok{stationary : III.6.5}
\NormalTok{triple, multiple : III.6.1}
\NormalTok{Sequences differing only in the order of their terms : III.6.1, R.7.8}
\NormalTok{Set : II.1.1}
\NormalTok{base : IV.1.3, 1.4}
\NormalTok{bounded, bounded above, bounded below : III.1.8, R.6.7}
\NormalTok{consisting of a single element : II.1.5}
\NormalTok{countable : III.6.4, R.7.7}
\NormalTok{decent : III.2.Ex.20}
\NormalTok{directed : III.1.10, R.6.8}
\NormalTok{empty : II.1.7, R.1.8}
\NormalTok{endowed with a structure : IV.1.4}
\NormalTok{equipotent to another set : R.7.1}
\NormalTok{finite : III.4.1}
\NormalTok{index (of a family) : II.3.4}
\NormalTok{inductive : III.2.4, R.6.9}
\NormalTok{infinite : III.6.1}
\NormalTok{left directed : III.1.10, R.6.8}
\NormalTok{Noetherian : III.6.5}
\NormalTok{of all x ∈ A such that P : II.1.6}
\NormalTok{of all x such that R (R a relation collectivizing in x) : II.1.4}
\NormalTok{of cardinals ≤ a : III.3.2}
\end{Highlighting}
\end{Shaded}

\begin{Shaded}
\begin{Highlighting}[]
\NormalTok{Set of classes of equivalent objects : II.5.9}
\NormalTok{of elements of a family : II.3.4, R.2.14}
\NormalTok{of mappings of E into F : II.5.2, R.2.2}
\NormalTok{of n elements (n an integer) : III.4.1}
\NormalTok{of objects of the form T for x ∈ A : II.1.6}
\NormalTok{of parameters of parametric representation : II.3.7}
\NormalTok{of σ{-}morphisms : IV.2.1}
\NormalTok{of subsets of a set : II.5.1, R.1.10}
\NormalTok{of subsets, of finite character : III.4.5}
\NormalTok{ordered : III.1.3, R.6.1}
\NormalTok{preordered : III.1.3, R.6.1}
\NormalTok{product : R.3.1, R.3.12, R.4.9}
\NormalTok{quotient : II.6.2, R.5.2}
\NormalTok{representative (of a relation) : II.3.1}
\NormalTok{right directed : III.1.10, R.6.8}
\NormalTok{totally ordered : III.1.12, R.6.4}
\NormalTok{transitive : III.2.Ex.20}
\NormalTok{universal : IV.3.1}
\NormalTok{well{-}ordered : III.2.1, R.6.5}
\NormalTok{whose only element is x : II.1.5}
\NormalTok{Sets, composition of : R.3.10}
\NormalTok{disjoint : II.4.7}
\NormalTok{equipotent : III.3.1}
\NormalTok{isomorphic : IV.1.5}
\NormalTok{theory of : II.1.1}
\NormalTok{Σ{-}admissible : IV.3.2}
\NormalTok{σ{-}morphism : IV.2.1}
\NormalTok{Sign : I.1.1, I.App.1}
\NormalTok{logical : I.1.1}
\NormalTok{rational : I.1.3}
\NormalTok{specific : I.1.1}
\NormalTok{substantific : I.1.3}
\NormalTok{Signed canonical extension : IV.2.Ex.1}
\NormalTok{Signed echelon type : IV.2.Ex.1}
\NormalTok{Significant sequence, word : I.App.2}
\NormalTok{Single{-}valued relation : I.5.3}
\NormalTok{Singular cardinal : III.6.Ex.17}
\NormalTok{Singular ordinal : III.6.Ex.16}
\NormalTok{Smaller than : IV.1.3}
\NormalTok{Solution (of an equation) : I.2.2}
\NormalTok{complete (or general) : I.5.2}
\NormalTok{of a universal mapping problem : IV.3.1}
\NormalTok{Source of a correspondence : II.3.1}
\NormalTok{Species, equivalent : IV.1.7}
\end{Highlighting}
\end{Shaded}

\begin{Shaded}
\begin{Highlighting}[]
\NormalTok{Species of algebraic structures : IV.1.4}
\NormalTok{of order structures : IV.1.4}
\NormalTok{of structures : IV.1.4, R.8.2}
\NormalTok{of topological structures : IV.1.4}
\NormalTok{poorer (richer) : IV.1.6}
\NormalTok{Specific sign : I.1.1}
\NormalTok{Stable subset : R.2.5}
\NormalTok{Stationary sequence : III.6.5}
\NormalTok{Strict upper bound : III.2.4}
\NormalTok{Strictly coarser (finer) structure : IV.2.2}
\NormalTok{Strictly decreasing (increasing) : III.1.5, R 6.12}
\NormalTok{Strictly greater (less) : III.1.3, R.6.3, R.7.3}
\NormalTok{Stronger theory : I.2.4}
\NormalTok{Strongly inaccessible cardinal : III.6.Ex.22}
\NormalTok{Structure, coarser : IV.2.2}
\NormalTok{deduced : IV.1.6}
\NormalTok{final : IV.2.5}
\NormalTok{finer : IV.2.2}
\NormalTok{generic : IV.1.4}
\NormalTok{induced : IV.2.4}
\NormalTok{initial : IV.2.3}
\NormalTok{of a set : IV.1.4, R.8.2}
\NormalTok{of species Σ : IV.1.4}
\NormalTok{order : R.6.1}
\NormalTok{product : IV.2.4}
\NormalTok{quotient : IV.2.6}
\NormalTok{richer : R.8.2}
\NormalTok{strictly coarser (finer) : IV.2.2}
\NormalTok{subordinate : IV.1.6}
\NormalTok{transport of : R.8.5}
\NormalTok{underlying : IV.1.6}
\NormalTok{Structures, comparable : IV.2.2}
\NormalTok{isomorphic : IV.1.5, R.8.6}
\NormalTok{Subfamily : II.3.5, R.2.14}
\NormalTok{Sublattice : III.4.Ex.9}
\NormalTok{Subordinate structure : IV.1.6}
\NormalTok{Subsequence (of a sequence) : III.6.1, R.8.7}
\NormalTok{Subset, cofinal : III.1.7, R.6.5}
\NormalTok{coinitial : III.1.7, R.6.5}
\NormalTok{consisting of a alone : R.1.9}
\NormalTok{empty : R.1.8}
\NormalTok{of a set : II.2.1, R.1.7}
\NormalTok{saturated (with respect to an equivalence relation) : II.6.4}
\NormalTok{symmetric : R.3.4}
\end{Highlighting}
\end{Shaded}

Subsets, disjoint : R.1.13 set of : II.5.1, R.1.10 stable : R.2.5 Substantific sign : I.1.3 Substitution, criterion of : I.1.2 Sum, cardinal : III.3.3 of a family of sets : II.4.8 of cardinals : III.3.3 of powers : R.7.5 of sets : R.4.5 ordinal : III.2.Ex.13 Supremum : III.1.9, R.6.7 Surjection : II.3.7, R.2.4 Surjective mapping : II.3.7, R.2.4 Symbol, functional : I.5.3 numerical : III.5.7 Symmetric graph : II.2.3 relation or subset : II.6.1, R.3.4 Symmetry, canonical : R.3.4 System, decimal : III.5.7 direct : III.7.5, R.6.13 dyadic : III.5.7 inverse : III.7.1, R.6.14 System of numeration : III.5.7 System of representatives : II.6.2

Target of a correspondence : II.3.1 Term : I.1.3 general : R.7.8 (first, kth, last) of a finite sequence : III.5.4 intrinsic : IV.1.6 nth : R.7.8 satisfying a relation : I.2.2 which can be put in the form T : I.5.2 Text, demonstrative : I.2.2 Theorem : I.2.2 of Cantor : III.3.6 of legitimation : I.3.3 of Zermelo : III.2.3, R.6.5 Theories, equivalent : I.2.4 Theory : I.1.1, I.2.1, I.2.2 contradictory : I.2.2 equalitarian : I.5.1 logical : I.3.1

\begin{Shaded}
\begin{Highlighting}[]
\NormalTok{Theory multivalent : R.8.7}
\NormalTok{of a species of structures : IV.1.4}
\NormalTok{of sets : II.1.1}
\NormalTok{of structures of a given species : R.8.2}
\NormalTok{quantified : I.4.2}
\NormalTok{stronger : I.2.4}
\NormalTok{univalent : R.8.7}
\NormalTok{Total order relation, total ordering : III.1.12}
\NormalTok{Totally ordered set : III.1.12, R.6.4}
\NormalTok{Trace of a family of sets : II.4.5}
\NormalTok{of a subset, or set of subsets : R.1.16}
\NormalTok{Transform of an element by a function : II.3.4, R.2.4}
\NormalTok{Transitive relation : II.6.1, R.5.1}
\NormalTok{Transitive set : III.2.Ex.20}
\NormalTok{Transitivity criteria : IV.2.3, 2.5}
\NormalTok{Transport of structure : R.8.5}
\NormalTok{Transportable relation : IV.1.3}
\NormalTok{Transporting a structure : IV.1.5}
\NormalTok{Transversal : II.6.2}
\NormalTok{Triple : II.2.2}
\NormalTok{Triple sequence : III.6.1}
\NormalTok{True relation : I.2.2}
\NormalTok{Typical characterization of a species of structures : IV.1.4}
\NormalTok{Typical quantifier : I.4.4}
\NormalTok{Typification : IV.1.3}

\NormalTok{Unbounded interval : III.1.13, R.6.4}
\NormalTok{Underlying structure : IV.1.6}
\NormalTok{Union, countable : R.7.9}
\NormalTok{of a family of sets : R.4.2}
\NormalTok{of a set of sets : II.4.1}
\NormalTok{of several sets : R.1.13}
\NormalTok{Univalent species of structures : IV.1.5}
\NormalTok{Univalent theory : R.8.7}
\NormalTok{Universal mapping : IV.3.1}
\NormalTok{Universal problem : IV.3.1}
\NormalTok{Universal quantifier : I.4.1}
\NormalTok{Universal set : IV.3.1}
\NormalTok{Upper bound : III.1.8, R.6.7}
\NormalTok{least : III.1.9, R.6.7}
\NormalTok{strict : III.2.4}

\NormalTok{Value, of a function : II.3.4}
\NormalTok{of a function at an element : R.2.1}
\end{Highlighting}
\end{Shaded}

Value, of a variable : R.1.2 taken by a correspondence : II.3.1 Variable : R.1.2 Variance of an assembly : IV.2.Ex.1

Weak compatibility (of an equivalence relation with a preorder relation): III.1.Ex.2

Zermelo's axiom (= axiom of choice) : R.4.9

Zermelo's theorem : III.2.3, R.6.5

Zorn's lemma : III.2.4, R.6.10

